\section{Introduction}
The computational demands of large language models (LLMs) are growing exponentially, driving training and inference costs to millions of dollars. High-performance GPU kernels are critical to reducing these costs. However, kernel development heavily relies on manual tuning by domain experts. As GPU architectures evolve rapidly and diversify, this manual optimization process has become increasingly complex and expensive, making automated GPU kernel generation a critical research direction for large-scale machine learning.

Recent advances in LLMs have demonstrated strong performance on general-purpose code generation tasks (e.g., Python, C++), making them promising tools for GPU kernel generation. 
\fktocheck{Existing approaches can be categorized into two classes: (1) specializing LLMs through fine-tuning or reinforcement learning (RL), and (2) building agent systems based on general-purpose LLMs.}

\fktocheck{Although RL has achieved success in decision-making tasks, its applicability to GPU kernel generation remains limited. High-quality kernel data are scarce, and kernel code must satisfy strict syntactic and semantic constraints. As a result, generated kernels often fail to compile or produce incorrect results. For example, AutoTriton reports an average correctness rate of less than 40\%.} 

% In contrast, agent approaches built on general-purpose LLMs avoid expensive training cost and can directly leverage pretrained models' strong code generation capabilities. Through multi-turn iterative refinement, these methods significantly improve correctness, offering better engineering feasibility and scalability. However, most existing agent systems rely on coarse-grained performance feedback to improve kernels and lack domain-specific kernel optimization knowledge. Without explicit guidance on what to optimize and how, agents often make random or conflicting modifications across iterations, resulting in limited performance gains. For example, GEAK shows little performance improvement after four optimization rounds.
In contrast, agent approaches built on general-purpose LLMs avoid expensive training costs and directly leverage pretrained models' strong code generation capabilities. 
\fktocheck{Through multi-turn iterative refinement, these methods significantly improve correctness, offering better engineering feasibility and scalability. However, most existing agent systems lack domain-specific kernel optimization knowledge and rely on coarse-grained performance feedback (e.g., execution time). This causes agents to make blind trial-and-error modifications, resulting in limited performance gains. For example, GEAK shows little improvement after four optimization rounds.}\fk{[first introduce GEAK and other methods, then summarize the weakness. Now the order is inversed.]}

The root cause lies in the intrinsic complexity of the kernel optimization space. Kernel performance is determined by intricate interactions among thread scheduling, memory hierarchy utilization, synchronization mechanisms, and hardware-specific features. Even minor code modifications can lead to orders-of-magnitude performance differences. Previous work has shown that even a small neural network subgraph on GPUs can exhibit an optimization space as large as $10^9$ configurations (Zhai et al., 2024). Without expert knowledge to systematically navigate this vast landscape, agents struggle to discover high-performance solutions.

% In contrast, agent approaches built on general-purpose LLMs avoid expensive training costs and directly leverage pretrained models' strong code generation capabilities. Through multi-turn iterative refinement, these methods significantly improve correctness, offering better engineering feasibility and scalability. However, most existing agent systems lack domain-specific kernel optimization knowledge and rely solely on coarse-grained performance feedback (e.g., execution time). This causes agents to make blind trial-and-error modifications, resulting in limited performance gains. For example, GEAK shows little improvement after four optimization rounds.

% This limitation comes from the intrinsic complexity of the kernel optimization space. Kernel performance is determined by intricate interactions among thread scheduling, memory hierarchy utilization, synchronization mechanisms, and hardware-specific features. Even minor code modifications can lead to orders-of-magnitude performance differences. Previous work has shown that even a small neural network subgraph on GPUs can exhibit an optimization space as large as $10^9$ configurations (Zhai et al., 2024). Without expert knowledge to systematically navigate this vast landscape, trial-and-error exploration struggles to discover high-performance solutions.

Based on these observations, we argue that fully unlocking the potential of LLMs for kernel generation requires integrating expert-level kernel optimization knowledge into agent systems to guide model decisions. To this end, we propose XX, an expert-guided agent framework that automatically generates high performance GPU kernels.

Specifically, we introduce expert-guided staged optimization, decomposing the kernel generation workflow into two hierarchical levels: high-level algorithmic structure design and fine-grained hardware-specific tuning. In the algorithmic structure design stage, we combine multi-seed search with algorithmic refinement to establish the performance upper bound. Multi-seed search generates diverse initial implementations to expand the exploration space. Algorithmic refinement then applies structural optimizations (e.g., operator fusion, algorithm replacement) to eliminate redundant computation and memory traffic at the source.
Subsequently, the hardware-specific tuning stage systematically realizes this performance potential through three sequential sub-stages with well-defined objectives: parallel mapping determines workload distribution across GPU cores, tensor tiling optimizes on-chip data reuse, and memory optimization refines data access patterns and pipeline execution. By decomposing the complex optimization into constrained sub-problems with explicit objectives, this staged approach effectively guides the LLM toward high-performance solutions. \fk{[overall ok, but need point out why expert guided?]}

% To effectively support this workflow, we introduce a stage-aware multi-agent collaborative framework consisting of a code agent, a profile agent, and a debug agent. Within each stage, the three agents collaborate on focused objectives through structured information exchange (e.g., bottleneck analysis, fix suggestions). Across stages, context propagation retains only finalized optimization decisions and filters out intermediate exploration details. This design prevents objective interference across stages and ensures cumulative improvements throughout the optimization process.
\fk{To support this workflow, we propose a stage-aware multi-agent collaboration mechanism consisting of a code agent, a profiling agent, and a debugging agent. Within each stage, the agents coordinate around clearly defined objectives via structured information exchange, including bottleneck identification and optimization proposals. Across stages, a selective context propagation strategy retains only finalized optimization decisions while discarding intermediate exploratory artifacts. This design reduces cross-stage objective interference and enables consistent, cumulative performance improvements over the course of the optimization process.}

The key contributions of this paper are as follows:
\begin{itemize}
\item We propose an expert-guided agent framework for GPU kernel generation that decomposes the optimization process into high-level algorithmic structure design and fine-grained hardware-specific tuning. By decomposing complex joint optimization into constrained sub-problems with explicit objectives, our approach effectively guides LLMs toward high-performance solutions.

\item We design a stage-aware multi-agent collaborative optimization mechanism, in which code, profile, and debug agents cooperate within each stage and propagate optimization decision across stages, achieving stable and accumulative improvements.

\item We systematically evaluate our approach on KernelBench and real-world workloads. Experimental results show that our approach achieves 100\% correctness while delivering up to $2.13\times$ speedup over PyTorch implementations, demonstrating the effectiveness of integrating expert knowledge into the kernel generation agent.
\end{itemize}